% Formato request/response OBD-II — tabla LaTeX
% Ejemplo con PID 0x0C (RPM). Requiere: array (ya en main.tex)

\begin{table}[!htbp]
\caption{Formato de petición y respuesta OBD-II (SAE J1979) encapsulados como payload ISO-TP sobre CAN. Se muestra el ejemplo de lectura de RPM (PID \texttt{0x0C}).}
\label{tab:obd2_format}
\centering
\renewcommand{\arraystretch}{1.4}

% --- Petición ---
\textbf{Petición} (Escáner \textrightarrow\ ECU, CAN ID: \texttt{0x7E0})

\vspace{0.3em}
\begin{tabular}{|c|c|c|c|c|c|c|c|c|}
\hline
& \textbf{PCI} & \textbf{Modo} & \textbf{PID} &
\textbf{Byte 3} & \textbf{Byte 4} & \textbf{Byte 5} & \textbf{Byte 6} & \textbf{Byte 7} \\
\hline
\textit{Valor} &
\texttt{0x02} & \texttt{0x01} & \texttt{0x0C} &
\textit{0xAA} & \textit{0xAA} & \textit{0xAA} & \textit{0xAA} & \textit{0xAA} \\
\hline
\textit{Significado} &
SF, 2 bytes & Modo 01 & RPM &
\multicolumn{5}{c|}{\textit{padding}} \\
\hline
\end{tabular}

\vspace{1em}

% --- Respuesta ---
\textbf{Respuesta} (ECU \textrightarrow\ Escáner, CAN ID: \texttt{0x7E8})

\vspace{0.3em}
\begin{tabular}{|c|c|c|c|c|c|c|c|c|}
\hline
& \textbf{PCI} & \textbf{Modo} & \textbf{PID} &
\textbf{Byte A} & \textbf{Byte B} & \textbf{Byte 5} & \textbf{Byte 6} & \textbf{Byte 7} \\
\hline
\textit{Valor} &
\texttt{0x04} & \texttt{0x41} & \texttt{0x0C} &
\texttt{0x1A} & \texttt{0xF0} &
\textit{0xAA} & \textit{0xAA} & \textit{0xAA} \\
\hline
\textit{Significado} &
SF, 4 bytes & Modo+\texttt{0x40} & RPM &
\multicolumn{2}{c|}{datos RPM} &
\multicolumn{3}{c|}{\textit{padding}} \\
\hline
\end{tabular}

\vspace{0.5em}
\begin{tabular}{ll}
\textbf{Decodificación:} & $\text{RPM} = \dfrac{(\texttt{A} \ll 8) \mid \texttt{B}}{4} = \dfrac{\texttt{0x1AF0}}{4} = 1724{,}0 \text{ rpm}$ \\
\end{tabular}
\end{table}
